% page 355 of the facsimile (image p-354 of the scan)
% block 170.2 x 198.8 mm ; left margin 31.8 mm
\pgc{10.160mm}{0.000mm}{
\l{}
\l{}
\l{}
\l{\cel{59}347}
\l{}
\l{\cel{4}\sou{lupulo}.-(\sou{bot}\cel{1}.) Planto ek la familio \textquotedbl{}urticei\textquotedbl{}, di qua la floro}
\l{\cel{7}esas uzata por facar biro. - L. humulus lupulus. - ISL.}
\l{}
\l{\cel{4}\sou{lupuso}. (\sou{patol.}) Pelo-morbo di naturo tuberklosala. - DEFIS.}
\l{}
\l{\cel{4}\sou{lurar}. (\sou{trans., per)}.I.Atraktar per ulo qua apetitigas.}
\l{\cel{7}II. (\sou{anke metaf.}) E.}
\l{}
\l{\cel{4}\sou{lustro}. I. Lumizilo suspendita, plura-brancha, uzata por lu-}
\l{\cel{7}mizar e dekorar la granda saloni, la kirki, la teatri.}
\l{\cel{7}- II. Periodo de kin yari. - DEFIRS.}
\l{}
\l{\cel{4}\sou{lustracar}.\cel{2}(\sou{trans.}) (en\cel{1}\sou{la epoki antiqua, en Roma)}. Purigar}
\l{\cel{7}lor ceremonio ta-skopa. - DEFIRS.}
\l{}
\l{\cel{4}\sou{lustrino}. Sorto di perkalino, tre apretita e briligita.- DEFIS.}
\l{}
\l{\cel{4}\sou{luto}. Induturo uzata por klozar hermetike vazi, por tegar}
\l{\cel{7}retorti, vitro-tubi o porcelano-tubi quin onu volas protek-}
\l{\cel{7}tat kontre la tro-varmeso di fairo. - EFIS.}
\l{}
\l{\cel{4}\sou{lutro}. (\sou{zool.}) Quadripedo karnavida, ek la familio \textquotedbl{}martri\textquotedbl{},}
\l{\cel{7}di qua la pilaro esas densa, qua vivas apud aquo, e qua su}
\l{\cel{7}nutras ek fishi. - L.\cel{1}\sou{lutra}. FIS.}
\l{}
\l{\cel{4}\sou{luxo}. I. Richeso, brilo, ostentata en la kozi di la vivo. -}
\l{\cel{7}II. (\sou{metaf.}) Kozo superflua. - DEFIS.}
\l{}
\l{\cel{4}luxacar. (\sou{trans}.)(\sou{pat.)}\cel{2}Igar (osto) ekirar de lua kavajo}
\l{\cel{7}normala. - DEFIS.}
\l{}
\l{\cel{4}luxurio. Peko kontre chasteso. - FIS.}
\l{}
\l{\cel{4}\sou{luzerno}. (\sou{bot.})\cel{2}Planto leguminosa quan onu kultivas en prati}
\l{\cel{7}artificala, kom nutrivo por la bestii. - L.\cel{1}\sou{medicago.}-}
\l{\cel{7}DEFIR.}
\l{}
\l{}
\l{}
\l{}
\l{\cel{31}---}
}
